\documentclass[10pt, a4paper]{article}
%
\usepackage{setspace}
\setstretch{1.2}      % [변경] 줄간격 1.2 -> 1.35로 확대하여 가독성 및 공간 확보
\usepackage{fontspec}
\usepackage{geometry}
% 여백 유지
\geometry{left=1.6cm, right=1.6cm, top=1.6cm, bottom=1.6cm}
\usepackage{titlesec}
\usepackage{enumitem}
\usepackage{hyperref}
\usepackage{kotex}
\usepackage{array}
\usepackage{xcolor}
\usepackage{fontawesome}

% 폰트 설정 (시스템에 Roboto 폰트가 설치되어 있어야 합니다)
\setmainfont{Roboto}[
    UprightFont = *-Regular,
    BoldFont = *-Bold,
    ItalicFont = *-Italic
]
\setsansfont{Roboto}

% [핵심] 얇고 작은 폰트를 위한 별도 정의 (Roboto Light 사용)
\newfontfamily\lightfont{Roboto Light}

% [변경] 본문 폰트 설정 수정
% 기존: 8pt / 10pt (너무 작음)
% 변경: 9.5pt / 13pt (가독성 확보 및 수직 공간 채움)
\newcommand{\descstyle}{\lightfont\fontsize{9.5pt}{10pt}\selectfont}

% [변경] 리스트 스타일 수정
% 기존: nolistsep (간격 없음)
% 변경: itemsep=0.3em (항목 간 숨쉴 공간 부여)
\setlist{leftmargin=1em, itemsep=0.3em} 

% 색상 정의
\definecolor{bonusSteelBlue}{RGB}{70, 130, 180}
\definecolor{darkGray}{RGB}{60, 60, 60}
\definecolor{lightGray}{RGB}{150, 150, 150}

% 하이퍼링크 설정
\hypersetup{
    colorlinks=true,
    linkcolor=bonusSteelBlue,
    urlcolor=bonusSteelBlue,
}

% 섹션 스타일
\titleformat{\section}
{\large\bfseries\color{darkGray}}
{}
{0em}
{}[\titlerule]

% 커스텀 명령어
\newcommand{\jobtitle}[2]{{\normalsize\textbf{#1}} \hfill {\descstyle\textit{#2}}}
\newcommand{\dates}[1]{\hfill {\descstyle\color{lightGray} #1}}

\begin{document}
\pagestyle{empty}

% --- HEADER START ---
\noindent
\begin{minipage}[t]{0.61 \textwidth} 
    {\huge\bfseries 두혁진} \\[0.5em]
    {\Large\textcolor{bonusSteelBlue}{system software engineer}} 
    \vspace{0.8em}
    
    {\descstyle 
    시스템의 기저 원리를 파고들며 운영 업무를 자동화하여 안정적인 서비스를 만드는 엔지니어입니다.
    C/C++ 기반의 견고한 시스템 프로그래밍 역량과 Rust의 메모리 안정성을 활용하여 
    소프트웨어의 신뢰성 확보에 기여하고 싶습니다.}
\end{minipage}%
\hfill
\begin{minipage}[t]{0.28\textwidth} 
    {\huge\bfseries contact} \\
    {\descstyle
    \faPhone \ 010-2805-0547 \\
    \faEnvelope \ \href{mailto:tyeirdoo06@gmail.com}{tyeirdoo06@gmail.com} \\
    \faGithub \ \href{https://github.com/hdoo42}{github.com/hdoo42}
    }
\end{minipage}

\vspace{1.5em} % [변경] 헤더와 본문 사이 간격 확대 (1.5 -> 2.0)

% --- 본문 시작 ---
\noindent
\begin{minipage}[t]{0.34\textwidth} % 왼쪽 컬럼

    % 1. EDUCATION
    \section*{\huge{education}}
    42 Seoul \\ 
    \dates{2021.11 - 2023.12}
    \vspace{0.3em}
    
    {\descstyle
    \begin{itemize}[label=\textbullet]
        \item C std lib, bash, IRC server 등을 직접 재구현하며 CS 기초 체화
        \item 엄격한 코드 규칙 준수 및 메모리 수동 관리를 통해 결함 없는 코드 작성 훈련
        \item 6대 coalition master(학생회장) 역임, 커뮤니티 공로 인정받아 IITP원장상 수상
    \end{itemize}
    }


    \vspace{1.5em} % [변경] 섹션 간 여백 확대

    % 2. SKILLS
    \section*{\huge{skills}}
   
    Programming \\
    {\descstyle Rust, C, C++, Bash, Python }
  
    \vspace{0.5em} % [변경] 항목 간 여백 확대
    
    Spoken Languages \\
    {\descstyle Japanese (Advanced Business) \\
    English (Daily / Simple Business) }      

    \vspace{0.5em} % [변경] 항목 간 여백 확대
    System \& OS \\
    {\descstyle Linux (Ubuntu), Shell Scripting }
    
    \vspace{0.5em} % [변경] 항목 간 여백 확대
    DevOps \& Tools \\
    {\descstyle Neovim, Git, Ansible, Docker, Prometheus, Grafana }

 
    \vspace{1.5em}

\end{minipage}
\hfill
\begin{minipage}[t]{0.62\textwidth} % 오른쪽 컬럼

    % 4. WORK EXPERIENCE
    \section*{\huge{work experience}}
    
    \jobtitle{(재)경산이노베이션아카데미}{System Software Engineer} \\
    \dates{2024.02 - 현재}
    \begin{itemize}[label=\textbullet] 
        \descstyle 
        \item Operational Automation (GsRoot)
        \begin{itemize}
            \item Rust/SQLite 기반의 워크스테이션 제어 서버 개발
            \item 관리자 부재 시간(새벽 등)의 대응 공백을 자동화로 해결하여, 기존 수 분~수 시간이 소요되던 장애 처리를 '즉시'로 단축 및 24시간 무중단 운영(Availability) 달성
        \end{itemize}
        
        \item System Integration (BLINK)
        \begin{itemize}
            \item 건물 출입 이벤트를 트리거로 LDAP 계정 상태(Unclose/Close)를 실시간 갱신하는 동기화 데몬 개발
            \item 교육생이 건물 내부 있는 경우에만 로그인 할 수 있도록 엑세스 권한 동기화
            \item 출석하지 않고 학습(로그인)시간 늘리는 부정행위 방지
        \end{itemize}
        
        \item Infrastructure Monitoring
        \begin{itemize}
            \item Prometheus/Slack 연동을 통해 기존의 수동 점검(Ansible ping)에 의존하던 장애 탐지를 자동화하여, 클러스터 상태 사각지대 제거
            \item 노드 다운 및 디스크 임계치(1GB 미만) 도달 시 alert 설정하여 실시간 현황 확인 및 대응
            
        \end{itemize}
        
        \item Automating Applicant Ops (Rust)
        \begin{itemize}
            \item 매일 반복되는 전일 지원자 데이터 수집 업무를 Rust로 자동화하여 수작업 비효율 제거, 기존 매일 수십분 걸리던 데이터 수집 및 갱신을 0시에 '미리' 완료
            \item 단순 크롤링을 넘어, 데이터 변경분(Upsert) 감지 및 정합성 처리를 통해 신뢰성 있는 데이터 파이프라인 구축
        \end{itemize}

        \item Technical Documentation
        \begin{itemize}
            \item Service Management Control 매뉴얼 작성 및 사내 지식 공유 위키(Wiki) 고도화
            \item Outline, n8n, NocoDB를 연동하여 사내 IT 이슈 처리 프로세스 자동화
            \item '이슈 등록 $\rightarrow$ 변경 $\rightarrow$ 문서화'의 워크플로우 정립을 통해 운영 노하우 자산화
        \end{itemize}
    \end{itemize}

    \vspace{1.2em} % [변경] 직장 간 여백 확대

    \jobtitle{(재)이노베이션 아카데미}{Intern} \\
    \dates{2023.10 - 2023.12}
    \begin{itemize}[label=\textbullet]
        \descstyle 
        \item Ansible을 활용한 클러스터 구성 관리
        \item 교육생 과제 명세 작성
    \end{itemize}

    \vspace{1.0em}

\end{minipage}

\end{document}
